\section{Performance Overhead}

While the primary goal of the additional coverage map is to increase path discovery, any added instrumentation inherently risks reducing the execution throughput. A significant drop in executions per second (exec/s) can offset the benefits of better feedback. This section evaluates the computational penalty introduced by our architecture.

\subsubsection{Execution Throughput}

To assess the impact, we compared the average executions per second across all campaigns for all tested configurations.

\begin{table}[htbp]
\centering
\begin{tabular}{lrrrrr}
\toprule
& \multicolumn{4}{c}{\textbf{Fuzzer Configuration}} \\
\cmidrule(l){2-5}
\textbf{Target} & \textbf{32 bit} & \textbf{64 bit} & \textbf{128 bit} & \textbf{Default} & \textbf{Overhead} \\
\midrule
\textbf{libxslt} & 586.28 & 644.88 & 587.78 & 698.71 & +13.22\% \\ 
\textbf{mruby} & 70.10 & 75.42 & 74.82 & 79.10 & +7.15\% \\ 
\textbf{openssl} & 2732.29 & 2610.04 & 2836.24 & 3644.51 & +25.2\% \\ 
\textbf{sqlite3} & 407.75 & 424.23 & 381.44 & 452.38 & +10.59\% \\ 
\textbf{systemd} & 115.13 & 110.06 & 111.76 & 110.53 & -1.62\% \\ 
\bottomrule
\end{tabular}
\caption{Average Execs Per Second and Resulting Overhead}
\label{tab:exec_overhead}
\end{table}

Table \ref{tab:exec_overhead} shows the average execution speeds and average overhead across all configurations. 

The data indicates that the secondary map, along with all the necessary mechanisms, introduces a measurable execution penalty that varies heavily by target. Across the benchmarks, the modified fuzzer experienced overhead ranging from $7\%$ to over $25\%$ --- with the anomaly of \texttt{systemd}, which can be attributed to environmental noise and statistical margin of error. 

Targets with particularly high densities of indirect branches exhibited the worst degradation. Furthermore, a significant proportion of the overhead resides in the post-execution phase, meaning that targets with higher absolute numbers of executions per second seem to suffer greater proportional degradation.

Some of the overhead is attributable to the injection of a callback, as the routine to call a function is computationally expensive. This could be mitigated in the future by injecting LLVM IR directly instead of relying on a function call.